\section{Conclusion}
Ce chapitre a permis de traduire le processus expérimental présenté au chapitre précédent en un système d’information structuré. Le choix de la méthode Merise a permis de représenter les acteurs, les flux d'informations et les enchaînements des traitements à travers le MCC, le MCT et le MOT. Cette modélisation a ainsi permis de conserver la logique du processus réellement observée sur le terrain, tout en identifiant les opérations pouvant être automatisées.

Sur cette base, l’application Agri Sens a été conçue sous la forme d’une progressive Web APP, en s’appuyant sur l’architecture FIWARE pour la gestion des données provenant du capteur 8-en-1 et les données provenant de la météo. L’application automatise notamment le calcul des indicateurs agronomiques, l'évaluation des besoins en fertilisation et en irrigation, ainsi que la visualisation des données et leur historique tout en conservant les interventions nécessitant une action humaine.

Ainsi, ce chapitre établit le passage entre le processus agronomique expérimenté et sa traduction informatique, en proposant une conception fidèle aux pratiques observées et aux besoins réels identifiés durant l'expérimentation.
